\section{Shared question and verified gap}
A robot cannot necessarily use computation that it wins: its queued actions may expire, or a moving object may leave the grasp tolerance before inference finishes. The original Compute Commons auction treated inference windows as interchangeable. We instead ask: \emph{Can an access policy allocate usable robot inference while preserving incentives to report private task values?}

\textbf{Economics:} what value and revenue must be traded for entrant access? \textbf{Computation:} how do robot-state constraints change feasible winners? \textbf{Behavior:} do operators distinguish winning compute from receiving a timely action? Their integration changes both the allocation rule and the meaning of access.

Real-time chunking addresses delayed action generation \cite{rtc}; FogROS2 addresses cloud robotics execution \cite{fogros}; deadline-aware task-and-motion planning already allocates reasoning effort \cite{sung}. Themis already auctions GPU resources \cite{themis}. Thus neither latency-aware robotics nor fair compute auctions are new. Our bounded contribution is an inspectable comparison of \emph{nominal versus usable access}, with robot-state ablations and critical payments on a shared feasible set. It is a synthetic mechanism demonstration, not a new control method or a hardware result.

\section{Strategic model and three lenses}
One committed coordinator schedules a batch of robot operators on one inference server. All jobs arrive at time zero. Operator $i$ privately values one timely result at $v_i\geq0$ and submits sealed bid $b_i\geq0$. Public telemetry fixes inference duration $\ell_i$, action-buffer steps $h_i$, control period $\Delta_i$, target speed bound $w_i$, pose tolerance $\epsilon_i$, network delay $\eta_i$ and guard $g_i$. In consistent time units, the usable deadline is
\begin{equation}\label{eq:deadline}
d_i=\min\{h_i\Delta_i,\epsilon_i/w_i\}-\eta_i-g_i.
\end{equation}
For $w_i=0$, the pose bound is infinite. This conservative drift surrogate assumes no predictive compensation. Local control and protective stopping remain outside the auction. Timeliness is necessary, not sufficient, for safe or successful manipulation.

The coordinator commits to horizon $T$, an end-of-horizon worker evaluation reserve $r$, verified entrant status $e_i\in\{0,1\}$ and credit $c\geq0$. Feasible sets $\mathcal F$ admit a nonpreemptive order with completion $C_i\leq d_i$ and $C_i\leq T-r$. Only value is strategically reported. The chosen set maximizes $\sum_{i\in S}(b_i+ce_i)$ over $S\in\mathcal F$, with a fixed tie order. Let $H_i^0$ be the highest rival-score sum over sets excluding $i$, and $H_i^1$ the highest rival-score sum over feasible sets containing $i$. Winners pay
\begin{equation}\label{eq:payment}
p_i=\max(0,H_i^0-H_i^1-ce_i);
\end{equation}
losers pay zero; individually infeasible jobs never win. One sealed deadline, allocation, execution and settlement end the round.

\textbf{Game theory:} utility is $v_ix_i-p_i$, where $x_i$ means timely service. A bid-independent feasible set makes winning monotone and (\ref{eq:payment}) its critical value. Truthful bidding is weakly dominant and hence a Bayesian Nash equilibrium for any private-value prior \cite{harsanyi,roughgarden}. This claim excludes strategic telemetry, status, arrivals, budgets and repeated rounds.

\textbf{Social choice:} firms seek task value, entrants usable access, workers evaluation authority, and the coordinator revenue. Credit is an explicit priority, not measured social benefit. We report task value, revenue, usable entrant access and reserved milliseconds separately. Sen's distinction between resources and opportunities motivates measuring usable access \cite{sen}.

\textbf{Mechanism design:} replace capacity-only admission by robot-state feasibility before charging critical payments. An EDF engineering baseline isolates the value of strategic allocation from the value of respecting deadlines.

\section{Computation and auction comparison}
Enumerate subsets, check each in earliest-deadline-first (EDF) order, maximize credited bids, then compute exclusion/inclusion thresholds. Complexity is $O(2^n n\log n+n^2 2^n)$ for this transparent implementation; it is an offline oracle, not a real-time scheduler. The capacity-only baseline uses the same optimization and EDF dispatch but omits robot deadlines from admission. Its truthful-input outputs are a diagnostic, not an equilibrium in the physical environment. EDF greedily admits feasible jobs by deadline and charges zero.

\begin{table}[h]
\caption{Actual synthetic batch: $T=400$ ms, $\ell_i=200$ ms, $v=(10,9,8,6)$ and $d=(200,200,400,400)$ for $A,B,D,E$. Only $E$ is eligible. Payments follow bidder order within each admitted set.}
\centering\small\setlength{\tabcolsep}{3pt}
\begin{tabular}{lccrr}\toprule
Rule & Admitted & Pays & Value & Rev.\\\midrule
Capacity only & $A,B$ & $8,8$ & 10 & 16\\
EDF & $A,D$ & $0,0$ & 18 & 0\\
Feasible & $A,D$ & $9,6$ & 18 & 15\\
Credit $c=3$ & $A,E$ & $9,5$ & 16 & 14\\
Reserve $r=100$ & $A$ & $9$ & 10 & 9\\
Both $r=100,c=5$ & $E$ & $5$ & 6 & 5\\\bottomrule
\end{tabular}\label{tab:worked}
\end{table}
Capacity-only admission charges $B$ for an expired result, yielding utility $-8$; bidding zero avoids that loss. The feasible rule restores the stated incentive domain. Credit exchanges two value units for one usable entrant result. Removing only buffer expiry or only pose freshness permits $A,B$ and value 19: both physical channels have a consequential role. Reservations reduce usable commercial time; their worker benefit remains unmeasured.

The code verifies 36,289 assertions, including independent permutation checks and unilateral deviations. A paired, seeded sweep of 300 synthetic six-robot batches compares all rules; Appendix A reports means and Monte Carlo intervals. These quantify the chosen input generator, not robot populations. Network-delay stress exposes the fragility of deadlines with no slack.

\section{Behavior and interdisciplinary integration}
The behavioral artifact asks for a bid, predicted payment, predicted timely completion and confidence before showing results. It records regret and a post-outcome explanation. A late baseline winner's underbid may be rational, whereas the same interpretation is unsupported for the verified feasible mechanism. This distinction changes the interface and the behavioral hypothesis; numeric bid error alone is inadequate \cite{hassidim}.

A planned counterbalanced pilot compares a payment tutorial with a physical-deadline explanation. Outcomes are payment error, timeliness prediction and utility regret, with legitimacy judgments measured separately. Learning and distrust remain competing explanations. Manual peer entry supports anonymous local export; it does not provide secure multiplayer bidding. No completed human interactions or classroom causal effects are claimed.

\section{Contribution, practical impact and roadmap}
The proposed case is a shared tabletop pick-and-place testbed used by small integrators and incumbent operators. A shared VLA inference service motivates the workload \cite{openvla}; local controllers handle protective stopping. Log queue depth, object motion, inference/network time and accepted action timestamps to test whether nominal access becomes executable service. Workers select evaluation tasks for the reserved period and can appeal unused capacity. Eligibility and telemetry need independent verification; payments cannot certify safety.

Academia needs calibrated traces and matched baselines; industry needs tail-latency and task-completion evidence; public institutions need usable access and an appeal route. Next: replay measured traces, then controlled robot trials and a preregistered comprehension pilot. Falsify the benefit if calibrated state constraints add no predictive or allocation value. Dynamic coupling, uncertain runtimes, solver latency and multiparameter incentives remain open.

\textbf{2056 aspiration.} Building on Vickrey, Harsanyi, Sen and Pearl \cite{vickrey,harsanyi,sen,pearl}, develop accountable allocation institutions whose incentives, physical feasibility and human effects are independently testable. International research recognition requires generalizable evidence beyond this course prototype.

\textbf{Open science and SDG 4.} Run \texttt{python} \path{code/robot_model.py}; outputs, notebook and source accompany this revision. The \href{https://github.com/dku-comsci-econ206-Autumn2026/-PS2-FP9-ComputeCommons}{course repository} and \href{https://huggingface.co/spaces/dku-comsci-econ206-2026/compute-commons}{Space} are the submission destinations. The intended SDG 4.4 contribution is a reproducible technical-learning exercise, evaluated through explanation accuracy; educational benefit has not been measured.
\IfFileExists{release.tex}{\noindent Computational snapshot: \href{https://github.com/Sunnyxinrong/-PS2-FP9-ComputeCommons/tree/46e71886f36476bd0712d26d4e41e2cfb86db213}{\texttt{46e7188}} (personal fork); \href{https://colab.research.google.com/github/Sunnyxinrong/-PS2-FP9-ComputeCommons/blob/46e71886f36476bd0712d26d4e41e2cfb86db213/code/compute_commons_ps2.ipynb}{runnable notebook}. Team integration and Space revision remain to be confirmed.
}{}
